% Exact characteristic-zero certificates for four determinant seven-planes.
% Depends on the necessary-prolongation lemma preceding R04/R-04.

\section{Four determinant seven-plane obstructions}

\begin{finding}{P-09}{two determinant types and their transposes have border rank at least fourteen}
In the rectangle numbering of the finite certificates, let
\[
 \begin{split}
 S_A&=\{0,1,3,6,7,9,12\},\\
 S_B&=\{0,1,3,6,7,12,13\}.
 \end{split}
\]
The corresponding determinant residual tensors have border
rank at least \(14\) over every algebraically closed field
of characteristic zero. The same holds for their transposes,
whose subsets are
\[
 S_A^{\mathsf t}=\{0,1,2,6,7,18,19\},\qquad
 S_B^{\mathsf t}=\{0,1,2,6,7,8,18\}.
\]
These are individual residual-orbit lower bounds; a uniform
seven-plane minimum also requires the other coordinate or
Borel-fixed cases.
\end{finding}
\begin{proof}
We supply the whole reduction and the finite exact certificate.
Let \(T\in A\otimes B\otimes C\) be one of the two tensors,
where \(\dim A=7\), and \(B,C\) are the spaces of its used
matrix variables. Their common dimension \(m\) is \(11\)
for type \(A\) and \(12\) for type \(B\). Every slice is a
nonsingular signed matching on its four-coordinate support,
and the union uses all \(m\) variables; therefore
\[
 U=T(C^*)\subset A\otimes B,\qquad \dim U=m.
\]

Use the necessary-prolongation condition proved before R-04.
If \(\BR(T)\leq13\), there is a thirteen-plane
\(L\supset U\) such that
\[
 \dim\ker\Psi_A(L)\geq13,\qquad
 \dim\ker\Psi_B(L)\geq13.
 \tag{*}
\]
Here \(\Psi_A:L\otimes A\to B\otimes\bigwedge^2A\) and
\(\Psi_B:L\otimes B\to A\otimes\bigwedge^2B\) are the
exterior maps. They differ from the earlier notation only
by factor permutations and an overall sign. In particular
\[
 \operatorname{rank}\Psi_A(L)\leq6\cdot13=78,\qquad
 \operatorname{rank}\Psi_B(L)\leq(m-1)13.
 \tag{**}
\]

Give a variable \(e_{ij}\) the row/column torus weight
\(R_i+C_j\), and give a slice label indexed by \((I,J)\)
the weight \(-\sum_{i\in I}R_i-\sum_{j\in J}C_j\).
This split torus stabilizes \(T\). It preserves \(U\), the
incidence \(L\supset U\), and both closed rank conditions
in (*). Choose an integral one-parameter subgroup separating the finitely
many ambient torus characters. Its limit in this projective incidence
is a direct sum of subspaces of the corresponding weight blocks and allows \(L\) to be torus-invariant. Hence \(L/U\) is a sum
of weight spaces.

Let \(V_w\) denote a full weight space of \(A\otimes B\).
The quotient \((A\otimes B)/U\) has \(47\) nonzero weight
blocks for type \(A\), and \(50\) for type \(B\).
These numbers and every weight coordinate are explicitly
enumerated in the certificate. For type \(B\),
\(\dim(L/U)=1\); its weight \(w\) satisfies
\(L\subset U+V_w\). For type \(A\),
\(\dim(L/U)=2\); either both dimensions lie in one weight
block of quotient dimension at least two, or they are
lines in two distinct blocks. There are \(14\) choices
of the former kind and \(\binom{47}{2}=1081\) of the
latter, for \(1095\) choices in total.

\paragraph{First stage: full weight-space extensions.}
For a chosen set of weights \(D\), put
\(L_D=U+\sum_{w\in D}V_w\). It contains every possible
\(L\) with those quotient weights. The certificate gives
a minor of one of the two exterior maps on \(L_D\)
of order
\[
 \dim(\text{its domain})-12
\]
with nonzero determinant modulo \(1000003\).
Consequently the corresponding kernel on \(L_D\) has
dimension at most \(12\), as does the kernel on any
subspace \(L\subset L_D\). This contradicts (*).
This stage excludes \(1042\) of the \(1095\) type \(A\)
choices, and \(47\) of the \(50\) type \(B\) choices.
In particular, no type \(A\) choice with both quotient
dimensions in one weight remains.

\paragraph{Second stage: a finite linear test on each weight line.}
We describe the test used on the remaining weights.
Fix \(w\) and write a quotient vector \(X\in V_w\)
in a basis of \(V_w/(U\cap V_w)\).
If \(U\cap V_w\ne0\), its nonzero vector has a unit
coefficient on every coordinate in its support, so one
coordinate may be omitted to choose this quotient basis.
If \(U\cap V_w=0\), use the whole coordinate basis.
Thus \(X\) has a nonzero parameter vector of dimension
at most five in all cases used below.

For an exterior map in a factor \(F\), let
\(\pi\) be the quotient of its target by the image of
that map on \(U\otimes F\). If \(f_1,\ldots,f_d\) are
the coordinate vectors of \(F\), then
\[
 \pi\Psi(X\otimes f_i)
\]
has weight \(w+\operatorname{wt}(f_i)\).
These weights are pairwise distinct. The rank increment
on adjoining \(X\) is therefore exactly the number of
nonzero projected columns.

The certificate supplies integer linear functionals that
annihilate \(\Psi(U\otimes F)\), separately in each weight
block. Their annihilation is checked as an exact integer
zero identity. Applying them to the new columns gives
matrices \(H_i\), linear in the quotient coordinates of
\(X\). If a projected column is zero, then \(H_iX=0\).
To prove a rank increment at least \(\ell\), it is enough
to check that for every set of \(d-\ell+1\) columns,
the stacked matrix of their \(H_i\)'s has full column
rank. This excludes a nonzero parameter vector for which
all those columns vanish. All these tests are finite
linear rank tests, certified modulo the same explicit
prime; no sampled or generic parameter is substituted.

\paragraph{Type \(B\).}
A certified minor gives
\[
 \operatorname{rank}\Psi_B(U)\geq137.
\]
The three remaining quotient weights have dimensions
\(3,5,3\). For each, the certificate verifies full column
rank for every one of the
\(\binom{12}{6}=924\) stacks corresponding to six zero
columns. Thus adjoining any nonzero line of one of these
weights increases rank by at least seven. This gives
\[
 \operatorname{rank}\Psi_B(L)\geq137+7=144
   >11\cdot13=143,
\]
contradicting (**).

\paragraph{Type \(A\).}
A certified minor gives
\[
 \operatorname{rank}\Psi_A(U)\geq72.
\]
The \(53\) remaining choices consist of pairs of distinct
weight lines. They use \(27\) quotient weight blocks.
For each such block \(w\), the finite linear test gives
a uniform increment lower bound \(\ell_w\), one of
\(2,3,4,6\), for the seven \(A\)-columns.
For two weights \(w,v\), let
\[
 F_w=\{w+\operatorname{wt}(a_s):1\leq s\leq7\},
 \qquad k_{w,v}=|F_w\cap F_v|.
\]
The projected column spaces of the two lines can intersect
only in these common weights; each line contributes at
most one vector in each weight. Their intersection
therefore has dimension at most \(k_{w,v}\).
The total rank increment is at least
\(\ell_w+\ell_v-k_{w,v}\).
For every one of the \(53\) remaining pairs, the certificate
checks
\[
 \ell_w+\ell_v-k_{w,v}>6.
\]
Consequently
\(\operatorname{rank}\Psi_A(L)\geq72+7=79>78\),
again contradicting (**).

All possible quotient weight choices have now been excluded.
Thus neither tensor has border rank at most \(13\).
Transposition is an invertible permutation of matrix
coordinates and preserves the determinant slices, proving
the same result for the two transposed subsets.
\end{proof}

\paragraph{Certificate and independent verifier.}
The complete data are in
\path{det7_two_type_certificate.json}.
The independent program
\path{verify_det7_certificate.py} reconstructs the
tensors, weights, and every required exterior-map entry
directly. It does not import the generator or use a
rank/nullspace oracle. It verifies the baseline minors,
all \(1089\) full-extension minors, all integer annihilator
identities, every required linear zero pattern, every
weight-overlap inequality, and complete coverage of the
\(1095+50\) quotient choices.
It uses integer NumPy arrays for modular arithmetic;
the prime is smaller than \(2^{20}\), all modular products
are smaller than \(2^{40}\), and explicit coefficient
bounds protect the unreduced integer zero identities.
The proof is therefore exact finite arithmetic, with no
floating-point tolerance. A nonzero minor modulo the
verified prime certifies a nonzero integer minor and
therefore the stated characteristic-zero result.

\paragraph{Geometric dictionary.}
The ambient construction is the proper Grassmannian
incidence with first-prolongation determinantal loci.
The torus limit gives a weight decomposition of a quotient
of the tautological bundle. Saturating selected weights
replaces a variable subbundle by an explicit larger
coordinate subbundle. The second stage tests linear
sections of projective quotient weight spaces; full-column
rank proves that the corresponding projective common zero
locus is empty. All maps and annihilator identities are
integral. Their nonzero modular minors certify the required
rank conditions in characteristic zero; no assertion about
all positive characteristics is made here.
